\subsection{Problem Framing}
\subsubsection{Sprint 1}
Game balancing is an essential consideration for game designers, to maximise the experience of all players. However, they are often posed with the issue of reliance on solely their own intuition and small-scale playtests to achieve this. These methods can introduce unconscious biases, and may not reflect the views of the game’s wider audience, for example users of varying playstyles. Subjective assessment makes it difficult for designers to identify problem areas where user engagement depreciates, such as difficulty spikes.

This project integrates an indie game with a telemetry application, generating intuitive live analytics for several actions as the game is run. The objective is to aid the designer’s understanding of a player’s experience, and the balance changes that could be applied to maximise the overall user experience.

\subsubsection{Sprint 2}
Game balancing is an essential consideration for the design and maintenance of games, as it directly impacts player engagement and overall user experience.
Game designers are often left to rely on solely their own intuition and small-scale playtests to navigate balancing decisions.
These methods can introduce unconscious biases, and their limitation in scope make it difficult to generalise solutions to a game's wider audience, such as users of varying play styles.

Consequently, designers lack objective and data-driven insights on how players engage with their game.
This presents a challenge in identifying problem areas, such as difficulty spikes and any other undesirable features where user experience is negatively impacted.

This project integrates an indie game with a telemetry application, generating intuitive analytics in real-time for several actions as the game is run.
The objective is to aid the designer's understanding of player behaviour, facilitating balance changes that are more informed to improve the overall user experience.
